% Part: methods
% Chapter: proofs
% Section: introduction

\documentclass[../../../include/open-logic-section]{subfiles}

\begin{document}

\olfileid{mth}{prf}{int}

\olsection{سريزه}

د منطق د لومړنيو درسونو د تجربې له مخې ښايي تاسو د
!!a{derivation} له نظام سره بلد ياست --- غالباً د طبيعي
استنتاج يا د فېچ په بڼه !!{derivation} نظام، يا د ثبوت
د ونې نظام۔ ښايي په ياد مو وي چې په دغو نظامونو کښې
مو ثبوتونه جوړول: يا مو !!a{formula} ثابتوله، يا مو
ښودله چې يو ټاکلے استدلال معتبر دے۔ د دې دپاره به مو
د نظام قاعدې تر هغې کارولې چې غوښتل شوې پايلې ته
ورسېدئ۔ د منطق \emph{په اړه} په استدلال کښې هم څيزونه
ثابتوو، خو په ډېرو حالاتو کښې د !!a{derivation} نظام
نه کاروو۔ په حقيقت کښې زموږ ډېر ثبوتونه په انګليسي
ژبه ليکل کېږي (کېدے شي ځينې سمبولونه هم ورسره وي)،
نه دا چې په بشپړه توګه د لومړۍ درجې منطق په ژبه وي۔
د داسې ثبوتونو په جوړولو کښې کېدے شي په لومړي سر
کښې حيران شئ: بې له !!a{derivation} نظامه يو څيز
څنګه ثابت کړم؟ له کومه ځايه پيل وکړم؟ څنګه پوه شم
چې زما ثبوت سم دے؟

د ثبوت له هڅې مخکې، دا پېژندل مهم دي چې ثبوت څه
دے او څنګه جوړېږي۔ لکه له نومه يې چې ښکاري،
\emph{ثبوت} د دې د ښودلو دپاره وي چې يو څيز رښتيا
دے۔ دا د خبرو اترو په بڼه هم تصورولے شئ: څوک
له تاسو پوښتنه کوي چې ايا کومه خبره رښتيا ده، د
بېلګې په توګه ايا له دوه پرته هر اولي عدد طاق دے۔
يوازې «هو» ويل بس نه دي؛ ښايي هغه وغواړي پوه شي
چې \emph{ولې}۔ په دې حالت کښې به ثبوت ورکوئ۔

په ورځنيو خبرو کښې ښايي ځواب ته لنډه اشاره يا
نيمګړے ځواب بس وي۔ خو په منطق او رياضياتو کښې
کره ثبوت غواړو: غواړو وښيو چې خبره له \emph{هر}
شک څخه ورهاخوا رښتيا ده۔ معنا يې دا ده چې د
ثبوت هر ګام بايد توجيه ولري، او توجيه يې بايد
سمه او قانع کوونکې وي (يعنې کارېدونکې فرضيه
بايد په هغه قضيه کښې په رښتيا فرض شوې وي چې
تاسو يې ثابتوئ؛ تعريفونه بايد سم وکارول شي؛
هغه استنتاجونه چې دليل پرې ولاړ وي بايد سم وي؛
او داسې نور)۔

عموماً موږ يوه وينا ثابتوو۔ دغو ويناوو ته بېلابېل
نومونه ورکوو: بيانونه، قضيې، لمونه، يا فرعي پايلې۔
بيان داسې بنسټيزه وينا ده چې ثبوت يې ارزښت لري:
د ثبتولو هومره مهمه وي، خو ښايي ډېره ژوره نه وي
يا ډېر ځله ونه کارول شي۔ قضيه يو مهم او د پام وړ
بيان دے۔ ثبوت يې ډېری وخت په څو پړاوونو وېشل کېږي،
او کله کله د هغه چا په نوم يادېږي چې لومړے يې
ثابته کړې وي (لکه د کانتور قضيه او د
لوېنهايم--سکولم قضيه)، يا د هغه حقيقت په نوم چې
اړيکه ورسره لري (لکه د بشپړتيا قضيه)۔ لم هغه
بيان يا قضيه ده چې د لا مهمې پايلې په ثبوت کښې
کارېږي۔ د مغشوشوونکې خبرې دا ده چې ځينې لمونه
پخپله هم مهمې پايلې وي او د وړاندې کوونکي په نوم
هم يادېږي (لکه د زورن لم)۔ فرعي پايله هغه نتيجه
ده چې له بلې پايلې په اسانۍ راځي۔

هغه وينا چې ثبوت يې غواړو، ډېر ځله فرضونه لري
چې څرګندوي د کومو څيزونو په اړه څه ثابتوو۔ ښايي
په دې پيل شي: «$!A$ دې د $!B \lif !C$ په بڼه
!!a{formula} وي»، يا «فرض کړئ $\Gamma \Proves !A$»،
يا ورته څه۔ دا د بيان، قضيې، يا لم \emph{فرضونه}
دي، او په خپل ثبوت کښې يې رښتيا ګڼلے شئ۔ دا
هغه څه محدودوي چې ثابتوو، او د بحث څيزونو ته
ځينې نومونه هم ورپېژني۔ د بېلګې په توګه، که ستاسو
بيان داسې پيل شي: «$!A$ دې د $!B \lif !C$ په بڼه
!!a{formula} وي»، نو تاسو يوازې د يوه ځانګړي ډول
فارمولونو (يعنې شرطي فارمولونو) په اړه څه ثابتوئ،
او څرګنده ده چې $!B \lif !C$ يو هر ډول شرطي
فارمول دے چې ستاسو ثبوت پرې خبرې کوي۔

\end{document}
